%
% vox-enable-disable.tex
%
% Z Specification for vox's per-repo enablement state machine.
% Models the three-state marker (Enabled / Dormant / Absent), the
% enable / disable / purge transitions, and the load-bearing
% biconditional  markerPresent <=> importPresent  of
% tool-enable-disable.md sec 2.11, together with the CLAUDE.md
% import-line writer idempotence of sec 2.4.
%
% Input design: docs/vox-enable-disable.md sec 4.
% Standard:     punt-kit/standards/tool-enable-disable.md sec 2.4, 2.7, 2.11.
%

\documentclass[a4paper,10pt,fleqn]{article}
\usepackage[margin=1in]{geometry}
\usepackage{fuzz}
\usepackage[colorlinks=true,linkcolor=blue,citecolor=blue,urlcolor=blue]{hyperref}

\begin{document}

\title{Vox Per-Repo Enablement: A Z Specification}
\author{Formal model of the \texttt{enabled} marker and the \texttt{enable} /
  \texttt{disable} verbs}
\date{July 2026}
\hypersetup{
  pdftitle={Vox Per-Repo Enablement: A Z Specification},
  pdfauthor={Punt Labs},
  pdfsubject={Z Specification},
  pdfcreator={fuzz/probcli}
}
\maketitle

\tableofcontents
\newpage

\section{Introduction}

A repository turns vox on with \texttt{vox enable} and off with
\texttt{vox disable}. Enablement is not a boolean field but a small state
machine over three observable facts: whether the tool-owned directory
\texttt{.punt-labs/vox/} exists, whether the \texttt{enabled} marker inside it
exists, and whether the repo's \texttt{CLAUDE.md} carries the canonical
\texttt{@.punt-labs/vox/CLAUDE.md} import. The standard
(\texttt{tool-enable-disable.md}) fixes one load-bearing property over these
three: the marker is present exactly when the import is present. This is the
\emph{enabled $\iff$ import} biconditional of \S2.11, and it is what
\texttt{punt audit} keys on and what every \texttt{enable}/\texttt{disable}
transition must preserve.

This specification pins the machine so that the invariant is a property of the
state schema rather than of scattered defensive checks, and so that the two
transitions the informal design underspecifies --- \texttt{Purge}, which the
\S4 state machine described only by its effect on the directory, and the
line-level idempotence of the import writer, which a single boolean cannot
express --- are stated precisely. Writing the model surfaced three places where
a literal reading of the design would strand an orphan import or resurrect a
directory; they are collected in Section~\ref{sec:findings}.

\section{Basic Types}

The model has no basic (given) types. Every fact it tracks is two-valued --- a
file or line is present or absent --- and the two counts it needs
(\texttt{importLines}) are bounded naturals. This is deliberate: the enablement
decision is entirely a question of presence, and a model that admits only
presence cannot smuggle in hidden state.

\section{Free Types}

A Z boolean, spelled to avoid the ProB/B keyword clash with \texttt{BOOL}. The
three facts of the state, and every derived flag, range over \texttt{ZBOOL}.

\begin{zed}
  ZBOOL ::= ztrue | zfalse
\end{zed}

\section{Global Constant}

\texttt{maxImports} bounds the number of canonical import lines the model
admits in a host \texttt{CLAUDE.md}. A conformant writer keeps this at most one;
the value $2$ is chosen so that an \emph{accidental duplicate} --- the state a
non-conformant or racing writer could leave, which \S2.4 requires
\texttt{disable} to collapse --- is representable and therefore checkable. Were
the bound $1$, the duplicate-collapse property would be vacuous, since the
offending state could not be written down.

\begin{axdef}
  maxImports : \nat_1
  \where
  maxImports = 2
\end{axdef}

\section{State}
\label{sec:state}

The enablement state is the three presence facts of the \S2.7 table:
\texttt{dirPresent} (the directory \texttt{.punt-labs/vox/} exists),
\texttt{markerPresent} (the \texttt{enabled} marker inside it exists), and
\texttt{importPresent} (the repo \texttt{CLAUDE.md} carries the canonical
import). The predicate carries the two invariants named in the design's \S4 and
in \S2.11 of the standard.

\begin{schema}{RepoEnablement}
  dirPresent : ZBOOL \\
  markerPresent : ZBOOL \\
  importPresent : ZBOOL
  \where
  markerPresent = importPresent \\
  markerPresent = ztrue \implies dirPresent = ztrue
\end{schema}

The first conjunct is the \S2.11 biconditional: over the two-valued
\texttt{ZBOOL}, $markerPresent = importPresent$ is exactly
$markerPresent = ztrue \iff importPresent = ztrue$. The second records that the
marker lives inside the directory, so a present marker forces a present
directory; its contrapositive, $dirPresent = zfalse \implies markerPresent =
zfalse$, is the fact that removing the directory removes the marker.

Only three of the eight assignments to the triple satisfy both conjuncts, and
they are exactly the three observable states of the \S2.7 table.

\begin{schema}{Enabled}
  RepoEnablement
  \where
  dirPresent = ztrue \\
  markerPresent = ztrue \\
  importPresent = ztrue
\end{schema}

\begin{schema}{Dormant}
  RepoEnablement
  \where
  dirPresent = ztrue \\
  markerPresent = zfalse \\
  importPresent = zfalse
\end{schema}

\begin{schema}{Absent}
  RepoEnablement
  \where
  dirPresent = zfalse \\
  markerPresent = zfalse \\
  importPresent = zfalse
\end{schema}

The states partition the reachable configurations: the biconditional forbids the
four assignments with $markerPresent \neq importPresent$, and $markerPresent =
ztrue \implies dirPresent = ztrue$ forbids the remaining outlier
($dirPresent = zfalse$, $markerPresent = ztrue$, $importPresent = ztrue$). What
is left is precisely \texttt{Enabled}, \texttt{Dormant}, \texttt{Absent} --- no
fourth state exists, which is the point of pinning the invariants in the schema.

\section{Initialisation}

A fresh repository --- one that has never run \texttt{enable} --- is
\texttt{Absent}: no directory, no marker, no import.

\begin{schema}{Init}
  RepoEnablement'
  \where
  dirPresent' = zfalse \\
  markerPresent' = zfalse \\
  importPresent' = zfalse
\end{schema}

\section{Transitions}
\label{sec:transitions}

Three operations move a repository between the states. Each preserves the
invariant by construction, because each fixes the primed triple to one of the
three legal assignments.

\subsection{Enable --- reach \texttt{Enabled} from anywhere}

\texttt{enable} deposits the guide, writes the marker, and adds the import if
absent (design \S3.3, standard \S2.3). It has no precondition: it reaches
\texttt{Enabled} from \texttt{Absent}, \texttt{Dormant}, or \texttt{Enabled}
alike, and from \texttt{Enabled} it changes nothing --- it is idempotent, which
is also the upgrade path.

\begin{schema}{Enable}
  \Delta RepoEnablement
  \where
  dirPresent' = ztrue \\
  markerPresent' = ztrue \\
  importPresent' = ztrue
\end{schema}

The postcondition is a constant, so $Enable \semi Enable$ has the same effect as
$Enable$: the second run reads an \texttt{Enabled} state and writes it back
unchanged. Idempotence at the \emph{line} level --- that the import appears
exactly once, never twice --- is not visible here, where \texttt{importPresent}
is a single boolean; it is the subject of Section~\ref{sec:writer}.

\subsection{Disable --- reach \texttt{Dormant} or \texttt{Absent}, non-destructively}

\texttt{disable} removes the import and the marker and deregisters settings, but
leaves the rest of the subtree in place (design \S3.3, standard \S2.9). The
directory is left exactly as it was found --- $dirPresent' = dirPresent$ --- so
\texttt{disable} lands in \texttt{Dormant} when the directory was present and in
\texttt{Absent} when it was not.

\begin{schema}{Disable}
  \Delta RepoEnablement
  \where
  dirPresent' = dirPresent \\
  markerPresent' = zfalse \\
  importPresent' = zfalse
\end{schema}

The frame $dirPresent' = dirPresent$ is the precise statement of
``non-destructive''. Writing the postcondition instead as $dirPresent' = ztrue$
--- ``leave the directory'' read as ``ensure a directory'' --- would resurrect a
directory when \texttt{disable} is run on an already-\texttt{Absent} repo,
manufacturing a spurious \texttt{Dormant} state from nothing
(finding~\ref{sec:findings}, item~2). The frame avoids that: \texttt{disable} on
\texttt{Absent} is a no-op that stays \texttt{Absent}.

\subsection{Purge --- reach \texttt{Absent} by removing the subtree}

\texttt{disable --purge} removes the whole subtree (design \S3.3). It is enabled
only when there is a directory to remove --- from \texttt{Dormant} or
\texttt{Enabled} --- and it lands in \texttt{Absent}.

\begin{schema}{Purge}
  \Delta RepoEnablement
  \where
  dirPresent = ztrue \\
  dirPresent' = zfalse \\
  markerPresent' = zfalse \\
  importPresent' = zfalse
\end{schema}

The postcondition sets \emph{all three} facts false, and the conjunct
$importPresent' = zfalse$ is load-bearing, not decorative. The import line lives
in \texttt{CLAUDE.md}, \emph{outside} the subtree \texttt{Purge} removes; a
\texttt{Purge} modelled as ``remove the directory'' alone --- $dirPresent' =
zfalse$, $markerPresent' = zfalse$, with \texttt{importPresent} left untouched
--- would, from \texttt{Enabled}, produce $(zfalse, zfalse, ztrue)$: a state
that violates the biconditional and leaves a dangling \texttt{@}-import that
\S2.11 flags as an orphan. \texttt{Purge} must therefore perform
\texttt{disable}'s import removal before deleting the subtree; it is
\texttt{Disable} followed by the directory removal, never the removal alone
(finding~\ref{sec:findings}, item~1).

\section{The Import-Line Writer}
\label{sec:writer}

The single boolean \texttt{importPresent} answers ``is the import there?'' but
not ``how many times?''. The \S2.4 idempotence contract is a statement about
\emph{lines}: appending when absent adds exactly one, a second \texttt{enable}
adds none, and \texttt{disable} removes every match, collapsing an accidental
duplicate to zero. To state and check those, the model refines the host file to
a line count.

\begin{schema}{ImportFile}
  importLines : \nat
  \where
  importLines \leq maxImports
\end{schema}

\texttt{importLines} is the number of canonical \texttt{@.punt-labs/vox/CLAUDE.md}
lines the host file carries, matched net of its terminator and outside code
blocks per \S2.4. It relates to the boolean of Section~\ref{sec:state} by
$importPresent = ztrue \iff importLines \geq 1$; the enablement state observes
only presence, while the writer observes the count.

\subsection{AppendImport --- append if absent, idempotently}

The \texttt{enable} step ``add the import if absent; never add it twice''
(\S2.3, \S2.4). From zero lines it writes one; from one or more it writes none.

\begin{schema}{AppendImport}
  \Delta ImportFile
  \where
  (importLines = 0 \implies importLines' = 1) \\
  (importLines \geq 1 \implies importLines' = importLines)
\end{schema}

$AppendImport \semi AppendImport$ from an empty file yields $0 \to 1 \to 1$: the
line appears exactly once, and the second \texttt{enable} adds nothing. Crucially
\texttt{AppendImport} \emph{never} raises the count above one on its own, so a
conformant writer cannot create a duplicate --- the $importLines = 2$ state is
reachable only through a non-conformant or racing writer, never through this
operation.

\subsection{RemoveImport --- remove every match}

The \texttt{disable} step removes the import. It removes \emph{every} matching
line, so a file that somehow carried two collapses to zero.

\begin{schema}{RemoveImport}
  \Delta ImportFile
  \where
  importLines' = 0
\end{schema}

This is the one operation that repairs a duplicate: from $importLines = 2$ it
reaches $0$. \texttt{AppendImport} preserves the audit's ``no duplicates''
property ($importLines \leq 1$) but does not by itself restore it from a
pre-existing duplicate --- only \texttt{RemoveImport}, i.e. \texttt{disable},
does (finding~\ref{sec:findings}, item~3).

\section{Predicate Schemas --- Invariants that are not Transitions}
\label{sec:predicates}

Three properties of the design are not transitions but constraints that must
hold in every state. They are stated as predicate schemas so that
\texttt{fuzz} type-checks them against the same state and \texttt{probcli} can
witness them.

\subsection{HookFires --- gating on the marker and the tool}

Vox's plugin session hooks fire in a repo only when the marker is present
\emph{and} the tool is installed (design \S3.5, standard \S2.7). The two-line
gate is exactly this conjunction; the \texttt{toolInstalled} conjunct is what
makes a clone of a marker-enabled repo on a box without vox a graceful no-op
rather than a ``command not found''.

\begin{schema}{HookFires}
  RepoEnablement \\
  toolInstalled : ZBOOL \\
  hookFires : ZBOOL
  \where
  hookFires = ztrue \iff \\
  \quad~ (markerPresent = ztrue \land toolInstalled = ztrue)
\end{schema}

A \texttt{Dormant} or \texttt{Absent} repo has $markerPresent = zfalse$, so
$hookFires = zfalse$ regardless of \texttt{toolInstalled}: vox is silent in a
repo that did not opt in. An \texttt{Enabled} repo on a box without the tool
($toolInstalled = zfalse$) is also silent --- the second conjunct fails --- which
is the \S2.11 ``no stale enabled tools'' behaviour at gate time.

\subsection{AuditPass --- the \S2.11 checks the model covers}

\texttt{punt audit} is clean when the biconditional holds, no import is an
orphan, and no import line is duplicated. The guide file that the import points
at is deposited into the directory by \texttt{enable} and left in place by
\texttt{disable}, so it is present exactly when the directory is:
$guideFilePresent = dirPresent$.

\begin{schema}{AuditPass}
  RepoEnablement \\
  ImportFile \\
  guideFilePresent : ZBOOL
  \where
  guideFilePresent = dirPresent \\
  importPresent = ztrue \iff importLines \geq 1 \\
  markerPresent = importPresent \\
  importPresent = ztrue \implies guideFilePresent = ztrue \\
  importLines \leq 1
\end{schema}

The ``no orphan import'' check --- $importPresent = ztrue \implies
guideFilePresent = ztrue$ --- is not an extra obligation but a consequence:
$importPresent = ztrue$ gives $markerPresent = ztrue$ (biconditional), which
gives $dirPresent = ztrue$ (state invariant), which gives $guideFilePresent =
ztrue$ ($guideFilePresent = dirPresent$). The only way to break it is a
\texttt{Purge} that fails to remove the import (finding~\ref{sec:findings},
item~1) --- exactly the orphan the audit exists to catch
(finding~\ref{sec:findings}, item~4).

\subsection{CliAgentGuided --- the intersection property}

The design's load-bearing intersection (\S1, \S3.4): after \texttt{enable} on a
box \emph{without} the plugin layer, the import is present and the deposited
guide carries the CLI-drive section, so a plugin-less agent can still drive vox
through the \texttt{vox} CLI. Because one static guide is deposited at both
scopes (\S2.5) and it carries the CLI section by construction, the CLI section is
present whenever the guide file is --- independent of the plugin.

\begin{schema}{CliAgentGuided}
  RepoEnablement \\
  pluginPresent : ZBOOL \\
  guideFilePresent : ZBOOL \\
  cliSection : ZBOOL
  \where
  guideFilePresent = dirPresent \\
  guideFilePresent = ztrue \implies cliSection = ztrue
\end{schema}

The intersection theorem specialises this to the plugin-less \texttt{enable}.
Compose \texttt{Enable} with \texttt{CliAgentGuided} read on the after-state and
fix $pluginPresent = zfalse$: \texttt{Enable} gives $importPresent' = ztrue$ and
$dirPresent' = ztrue$; the latter forces $guideFilePresent = ztrue$
($guideFilePresent = dirPresent$), which forces $cliSection = ztrue$. So after a
plugin-less \texttt{enable} both $importPresent$ and $cliSection$ hold --- a
working CLI-driving agent, the design's explicit acceptance scenario. The
plugin-less box is not a special case in the model; it is the general case read
at $pluginPresent = zfalse$, and the CLI section is present there for the same
reason it is present everywhere (finding~\ref{sec:findings}, item~5). The theorem
is stated as a composition rather than a new operation deliberately: folding it
into an \texttt{Enable}-plus-observation schema would make \texttt{CliAgentGuided}
an auxiliary of that schema and cost \texttt{Enable} its own standing as a
transition to animate.

\section{Key Properties}

The following hold by construction (state invariants) or as operation
postconditions; they are the intended proof obligations for animation.

\begin{itemize}
  \item \textbf{The biconditional is universal}: $markerPresent = importPresent$
    is a state invariant, so every reachable state satisfies it. Each of the
    three transitions writes one of the three legal triples, so none can leave a
    state that violates it. This is the property model-checked in
    Section~\ref{sec:modelcheck}.
  \item \textbf{Exactly three states}: the two invariants admit only
    \texttt{Enabled}, \texttt{Dormant}, \texttt{Absent}; the biconditional rules
    out the four mismatched triples and $markerPresent \implies dirPresent$ rules
    out the last outlier.
  \item \textbf{Enable is total and idempotent}: \texttt{Enable} has no
    precondition, reaches \texttt{Enabled} from any state, and $Enable \semi
    Enable = Enable$; at the line level $AppendImport \semi AppendImport$ takes
    $0 \to 1 \to 1$, so the import appears exactly once.
  \item \textbf{Disable is non-destructive}: $dirPresent' = dirPresent$, so the
    directory survives; \texttt{disable} reaches \texttt{Dormant} from
    \texttt{Enabled} and stays \texttt{Absent} from \texttt{Absent}, never
    creating a directory.
  \item \textbf{Disable collapses a duplicate}: \texttt{RemoveImport} sets
    $importLines' = 0$ from any count, so a $importLines = 2$ file is repaired to
    $0$ and the audit's $importLines \leq 1$ is restored.
  \item \textbf{Purge clears the import too}: \texttt{Purge} sets
    $importPresent' = zfalse$ as well as $dirPresent' = zfalse$, so it leaves no
    orphan import; a directory-only purge would violate the biconditional.
  \item \textbf{No orphan import}: $importPresent \implies guideFilePresent$ is a
    theorem given $guideFilePresent = dirPresent$ and the state invariants, not a
    separate runtime check.
  \item \textbf{Hooks gate on marker and tool}: $hookFires \iff (markerPresent
    \land toolInstalled)$, so vox is silent in a non-enabled repo and a graceful
    no-op on a box without the tool.
  \item \textbf{CLI guidance is plugin-independent}: after \texttt{enable} the
    CLI section is present whenever the guide file is, so a plugin-less box that
    enables yields an agent that can drive vox through the CLI.
\end{itemize}

\section{Model Checking}
\label{sec:modelcheck}

The biconditional is the property worth exploring exhaustively, because it is
the one \texttt{punt audit} keys on and the one a wrong transition could break
silently. The state is tiny --- three booleans, hence eight raw configurations,
of which the invariant admits three --- so the reachable state space of
\texttt{Init} together with \texttt{Enable}, \texttt{Disable}, and \texttt{Purge}
is finite and small.

\texttt{probcli} explores it in full: from \texttt{Init} ($=$ \texttt{Absent}),
\texttt{Enable} reaches \texttt{Enabled}, \texttt{Disable} from \texttt{Enabled}
reaches \texttt{Dormant}, \texttt{Disable} from \texttt{Dormant} or
\texttt{Absent} stays there, \texttt{Purge} from \texttt{Enabled} or
\texttt{Dormant} reaches \texttt{Absent}. Three states are reachable ---
\texttt{Absent}, \texttt{Dormant}, \texttt{Enabled} --- and in every one
$markerPresent = importPresent$ holds. No sequence of \texttt{Enable},
\texttt{Disable}, \texttt{Purge} reaches a state in which the marker and the
import disagree; the invariant is never violated and the machine does not
deadlock.

\section{Findings from Formalisation}
\label{sec:findings}

Writing the model surfaced these decisions, which the implementation design must
realise consistently. Each is a place where the informal \S4 state machine was
silent, underspecified, or would --- read literally --- break the biconditional.

\begin{enumerate}
  \item \textbf{\texttt{Purge} must remove the import line, not only the
    directory.} The design's \S4 lists \texttt{Purge}'s postcondition as
    $\neg dirPresent$ alone. But the import line lives in \texttt{CLAUDE.md},
    outside the subtree \texttt{Purge} deletes, while the marker lives inside it.
    A \texttt{Purge} that removes only the subtree would, from \texttt{Enabled},
    leave $(dirPresent, markerPresent, importPresent) = (zfalse, zfalse, ztrue)$
    --- a state that both violates the \S2.11 biconditional and leaves a dangling
    \texttt{@}-import that 404s at read time (the \S2.11 ``no orphan imports''
    hard fail). The model therefore makes $importPresent' = zfalse$ a required
    postcondition of \texttt{Purge}: \texttt{disable --purge} is \texttt{disable}
    (which removes the import) \emph{then} the subtree removal, never the subtree
    removal alone. The prose of \S3.3 sequences it correctly; the \S4 state
    machine's one-line summary does not, and an implementation that follows only
    \S4 would ship the orphan.

  \item \textbf{\texttt{Disable}'s non-destructiveness is $dirPresent' =
    dirPresent$, not $dirPresent' = ztrue$.} ``Leave the directory in place'' is
    a frame, not an assertion of presence. If \texttt{disable} instead
    \emph{ensured} a directory, running it on an already-\texttt{Absent} repo
    would create an empty \texttt{.punt-labs/vox/} with no marker --- a spurious
    \texttt{Dormant} state conjured from nothing, and a directory the tool did
    not just create, against the org ``never delete/create what you did not
    handle'' rule. The frame $dirPresent' = dirPresent$ is the correct reading:
    \texttt{disable} reaches \texttt{Dormant} only from a state that already had
    the directory.

  \item \textbf{The boolean \texttt{importPresent} cannot express the \S2.4
    duplicate-collapse; the line count is required.} The biconditional treats
    ``one import'' and ``two imports'' alike --- both are $importPresent =
    ztrue$ --- so ``\texttt{disable} removes every match'' and ``a second
    \texttt{enable} adds none'' are invisible at the boolean level. The
    \texttt{ImportFile} machine ($importLines : 0 \upto maxImports$, with
    $maxImports = 2$) is what makes \texttt{AppendImport}'s idempotence and
    \texttt{RemoveImport}'s total removal checkable. \texttt{AppendImport}
    \emph{preserves} the audit's $importLines \leq 1$ but does not by itself
    repair a pre-existing duplicate; only \texttt{RemoveImport} --- i.e.
    \texttt{disable} --- collapses $2$ to $0$. A design that models the import as
    a boolean would silently drop the requirement that \texttt{disable} heal a
    racing writer's duplicate.

  \item \textbf{``No orphan import'' is a theorem, not an independent check.}
    Once the deposited guide is present exactly when the directory is
    ($guideFilePresent = dirPresent$ --- \texttt{enable} deposits it,
    \texttt{disable} leaves it dormant), the \S2.11 orphan check $importPresent
    \implies guideFilePresent$ follows from the two state invariants. The single
    way to falsify it is the directory-only \texttt{Purge} of item~1. So the
    orphan the audit forbids and the \texttt{Purge} bug are the same defect seen
    from two sides, and fixing item~1 discharges the orphan check for free.

  \item \textbf{The intersection property is unconditional on the plugin.} The
    ``a plugin-less box that enables yields a CLI-driving agent'' scenario is not
    a special case: because one static guide is deposited at both scopes (\S2.5)
    and it carries the CLI-drive section by construction, $cliSection$ tracks
    $guideFilePresent$, which tracks $dirPresent$, with no dependence on
    $pluginPresent$. \texttt{CliAgentGuided} holds at $pluginPresent = zfalse$
    for the same reason it holds everywhere. The design implication is a
    constraint on the guide's \emph{content}: its CLI section must not be gated
    on plugin absence --- it is always present, and the preamble merely tells the
    agent which surface to prefer. A guide that emitted the CLI section only when
    the plugin was absent would still satisfy the scenario but would break the
    single-static-file contract of \S2.5.
\end{enumerate}

\end{document}
